\section{Introduction}\label{sec:introduction}

Benchmark scores drive model selection and research priorities for large language models, yet training data contamination is pervasive: trillion-token web crawls routinely ingest benchmark items \citep{oren2024proving}, and practitioners cannot determine how much of a reported score reflects genuine capability versus memorization \citep{hendrycks2021mmlu,chang2025benchmarking}. Despite twenty mitigation strategies \citep{sun2025emperor}, contamination-free benchmarks \citep{white2024livebench}, and dynamic regeneration \citep{dong_2024_cdd}, a fundamental question remains: what is the actual effect of contamination on aggregate scores? The implicit assumption is uniform inflation, but this has never been tested across model families.

We perform controlled SFT injection (LoRA, five dosage levels) into six models spanning five independent families from five labs (Alibaba, Microsoft, Allen AI, Google, TII) and discover a spectrum of dose-response patterns (Figure~\ref{fig:dose_response})---from immunity (Phi-4, $\pm$0.5pp) and benefit ($+$3--4pp, three families) to collapse (Qwen-14B, $-$9 to $-$25pp across seeds).

\begin{figure}[t]
\centering
\includegraphics[width=\columnwidth]{figures/fig_dose_response.pdf}
\caption{MMLU accuracy vs.\ contamination dosage (5-shot, seed$=$42; LLaMA-3.1 at seed$=$0). Immune and Benefit replicate across families; Collapse and V-recovery are preliminary (Qwen2.5-specific). Cross-seed stability in Table~\ref{tab:cross_seed}.}
\label{fig:dose_response}
\end{figure}

Across six models, three lower-baseline models benefit ($+$3.2 to $+$4.2pp), one is immune (Phi-4), and the remaining two (both Qwen2.5) show harm---with the 14B model 5$\times$ more fragile than 7B. These effects are benchmark-conditional: the same model can collapse on MMLU ($-$9.4pp) yet gain on GSM8K ($+$13pp). ContamCal, our diagnostic framework combining multi-signal exposure mapping with contamination-augmented IRT (2PL+$\gamma_k$), enables this characterization.

Our contributions:
\begin{enumerate}
    \item \textbf{Dose-response profiles with cross-seed and cross-benchmark validation.} Controlled injection across five families reveals two replicated patterns (Immune, Benefit) and two preliminary Qwen2.5-specific patterns (Collapse, V-recovery). Out-of-sample validation on LLaMA-3.1-8B confirms generalizability. Cross-benchmark transfer reveals directional reversal on GSM8K for three families. Multi-seed evaluation ($n\!=\!3$) shows direction is reproducible but magnitude is seed-sensitive.
    \item \textbf{Domain-level contamination diagnosis via IRT.} ContamCal embeds multi-signal RRF exposure scores in a contamination-augmented IRT model (2PL+$\gamma$)---an application of standard explanatory IRT \citep{deboeck_2004_explanatory_irt} to LLM contamination---enabling domain-level diagnosis with credible intervals. Domain analysis reveals opposing contamination directions (Social Sciences $\gamma_d < 0$, STEM $\gamma_d > 0$) invisible to aggregate scoring. No item-specific memorization is detected via difference-in-differences (MDES 1.25--2.05pp).
    \item \textbf{Diagnostic pipeline requirements and signal failure.} Min-K\%++ \citep{zhang_2024_minkpp} exhibits directional reversal at 0.5B (paired Cohen's $d = -1.98$ at d100). RRF-based multi-signal fusion is essential: a single-signal pipeline produces severe model misfit (PPP $< 0.001$). Simulation establishes $\gamma$ recovery requires $J \geq 50$ at AUC $\geq 0.55$, above current practice ($J = 8$--$30$).
\end{enumerate}

\noindent Code and data are available at \url{https://github.com/anonymous-nlp-2026-2/contamcal-irt-calibration}.
